\documentclass[conference]{IEEEtran}
\IEEEoverridecommandlockouts
\usepackage{cite}
\usepackage{amsmath,amssymb,amsfonts}
\usepackage{algorithmic}
\usepackage{graphicx}
\usepackage{textcomp}
\usepackage{xcolor}
\usepackage{booktabs}
\usepackage{url}
\def\BibTeX{{\rm B\kern-.05em{\sc i\kern-.025em b}\kern-.08em
    T\kern-.1667em\lower.7ex\hbox{E}\kern-.125emX}}

\begin{document}

\title{TactStyle: Enhancing 3D Model Stylization with Realistic\\
Texture Generation for Digital Fabrication}

\author{
\IEEEauthorblockN{Omprakash Pugazhendhi}
\IEEEauthorblockA{
  \textit{Department of Computer Science and Engineering}\\
  \textit{Vellore Institute of Technology}\\
  Chennai, India\\
  omprakash.2021@vitstudent.ac.in
}
}

\maketitle

\begin{abstract}
Existing 3D model stylization pipelines transfer the visual appearance of a reference
texture onto a target surface but neglect the tactile dimension that is critical for
digital fabrication applications such as 3D printing. A surface stylized with a
wood-grain image may look accurate but feel completely smooth when printed because the
geometry is never modified. We present a unified machine learning system that addresses
both visual and tactile fidelity simultaneously. Our pipeline combines neural style
transfer methods with a fine-tuned diffusion model that generates displacement
heightfields from single reference texture images. The heightfields are then applied
as vertex displacements to arbitrary 3D meshes via UV parameterization. We train the
system on the Describable Textures Dataset containing 5640 images across 47 categories
and evaluate generated heightfields using Structural Similarity Index, Learned
Perceptual Image Patch Similarity, and Frechet Inception Distance. Our diffusion-based
approach achieves an SSIM of 0.387, LPIPS of 0.334, and FID of 51.2 on the test set,
outperforming gradient-based synthetic baselines. We also introduce texture interpolation
in pixel, feature, and CLIP embedding spaces, and release an interactive Gradio
demonstration that produces a downloadable textured 3D model from a single image upload.
\end{abstract}

\begin{IEEEkeywords}
texture synthesis, neural style transfer, heightfield generation, diffusion models,
digital fabrication, 3D stylization, tactile rendering
\end{IEEEkeywords}

%% ─────────────────────────────────────────────────────────────────────────────
\section{Introduction}
%% ─────────────────────────────────────────────────────────────────────────────

Digital fabrication techniques such as fused deposition modeling and
stereolithography have made it possible for non-experts to produce complex
physical objects at low cost. A natural next step is to allow users to
specify not only the shape of a printed object but also its surface character:
a smartphone case that feels like woven fabric, a prosthetic grip that
replicates the texture of cork, or a tactile map whose rivers and roads can
be distinguished by touch alone.

Current stylization workflows treat a texture as a 2D color image that is
mapped onto a surface without modifying the underlying geometry. This approach
produces visually convincing renders but physically flat objects. To transfer
both the visual and tactile character of a reference texture, the surface
geometry must be perturbed to encode the micro-relief of the material.

Heightfields — grayscale images in which each pixel encodes a surface
displacement along the local normal — provide a compact and mesh-agnostic
representation of surface texture. Prior work on tactile texture generation,
notably TactStyle \cite{b1}, demonstrated that a modified diffusion model
conditioned on a reference image can generate plausible heightfields and that
3D-printed objects stylized with these heightfields are perceptually
indistinguishable from objects printed with ground-truth scan data in
psychophysical studies.

In this paper I present a complete open-source system that:
\begin{itemize}
  \item classifies input textures using ResNet-50 \cite{b2} fine-tuned on the
        Describable Textures Dataset (DTD) \cite{b3};
  \item synthesizes stylized images using three neural style transfer methods
        — Gram-matrix optimization \cite{b4}, fast residual networks \cite{b5},
        and Adaptive Instance Normalization (AdaIN) \cite{b6};
  \item generates displacement heightfields via Stable Diffusion 1.5
        \cite{b7} fine-tuned with Low-Rank Adaptation (LoRA) \cite{b8} on
        synthetic texture-heightfield pairs derived from DTD;
  \item applies heightfields as vertex displacements to arbitrary 3D meshes
        using UV parameterization with xatlas \cite{b9};
  \item supports texture interpolation in pixel space, AdaIN feature space,
        and CLIP \cite{b10} embedding space; and
  \item provides an interactive Gradio \cite{b11} demonstration.
\end{itemize}

I also integrate a bridge to 3DShape2VecSet \cite{b12}, a transformer-based
diffusion model for 3D shape generation, allowing meshes generated from image
prompts to be textured by the same pipeline.

%% ─────────────────────────────────────────────────────────────────────────────
\section{Related Work}
%% ─────────────────────────────────────────────────────────────────────────────

\subsection{Neural Style Transfer}

Gatys et al.~\cite{b4} introduced the seminal approach of separating content
and style in the feature space of a pre-trained VGG-19 convolutional network.
Style is encoded as the Gram matrix of activations at multiple layers;
optimization over the input image minimizes a weighted sum of content loss
and style loss. While this method produces high-quality results, it requires
hundreds of forward-backward passes and is prohibitively slow for interactive
use.

Johnson et al.~\cite{b5} addressed speed by training a feed-forward residual
network on a fixed style image, enabling real-time stylization of arbitrary
content images at inference time. Huang and Belongie~\cite{b6} generalized
this to arbitrary styles at inference time through Adaptive Instance
Normalization, which aligns the channel-wise mean and variance of content
features to those of style features before decoding.

\subsection{Tactile Texture Generation}

TactStyle \cite{b1} first demonstrated end-to-end tactile texture generation
from a single image using a modified image-to-image diffusion model.
The system applies heightfields to 3D models via UV mapping and a Blender
plugin interface. A psychophysical study validated that participants could not
reliably distinguish TactStyle-generated textures from direct scan-based
fabrication. Our work extends this approach with open-source training code,
multiple NST baselines, texture interpolation, CLIP conditioning, and an
evaluation suite.

\subsection{Language-Guided Haptic Synthesis}

Recent work \cite{b13} proposes a bimodal VAE with a CLIP-aligned latent
space encoding both sliding vibration models and tapping transient banks, with
a decoupled Stable Diffusion 2 branch for visual textures. This demonstrates
that CLIP embeddings provide a natural bridge between language, vision, and
touch. We adopt the CLIP image encoder as a conditioning signal for texture
interpolation in our system.

\subsection{3D Shape Representation}

3DShape2VecSet \cite{b12} encodes 3D shapes as sets of latent vectors
processed by a transformer diffusion model, supporting image-conditioned
shape generation. Our pipeline is compatible with meshes produced by this
system, enabling an end-to-end flow from a reference image to a generated
3D shape with an applied tactile texture.

%% ─────────────────────────────────────────────────────────────────────────────
\section{System Architecture}
%% ─────────────────────────────────────────────────────────────────────────────

The system consists of five sequential modules as shown in Fig.~\ref{fig:arch}.

\begin{figure}[htbp]
\centerline{\includegraphics[width=0.48\textwidth]{figures/architecture.pdf}}
\caption{System architecture. A reference image passes through the texture
         classifier, then optionally through neural style transfer, before
         the heightmap generator produces a displacement field that is
         applied to a 3D mesh via UV parameterization.}
\label{fig:arch}
\end{figure}

\subsection{Texture Classifier}

I fine-tune a ResNet-50 backbone \cite{b2} pre-trained on ImageNet \cite{b14}
on all 5640 images of DTD. Layers \texttt{conv1} through \texttt{layer2} are
frozen; \texttt{layer3}, \texttt{layer4}, and a newly added 47-class linear
head are trained end-to-end. AdamW with learning rate $10^{-4}$, cosine
annealing over 30 epochs, and label smoothing of 0.1 are used. The classifier
is used both for dataset organization and as an optional pre-filter to select
the most relevant training subset for a given style image.

\subsection{Neural Style Transfer}

Three NST methods are implemented and compared:

\textbf{Gatys (VGG-19).} Content features are extracted at \texttt{relu4\_2};
style features are extracted at \texttt{\{relu1\_1, relu2\_1, relu3\_1,
relu4\_1, relu5\_1\}}. The total loss is
\begin{equation}
  \mathcal{L} = \alpha\,\mathcal{L}_c + \beta\,\mathcal{L}_s + \gamma\,\mathcal{L}_{tv},
  \label{eq:gatys}
\end{equation}
where $\mathcal{L}_c$ is the mean-squared content feature distance,
$\mathcal{L}_s$ is the sum of Gram-matrix distances across style layers,
and $\mathcal{L}_{tv}$ is a total-variation regularizer. LBFGS is used
as the optimizer.

\textbf{Johnson (Fast).} A residual encoder-decoder with instance
normalization is trained once per style image using perceptual losses
computed on VGG-16 features. Inference is a single forward pass.

\textbf{AdaIN.} Given content features $\mathbf{x}$ and style features
$\mathbf{y}$, the transfer is computed as
\begin{equation}
  \mathrm{AdaIN}(\mathbf{x}, \mathbf{y}) =
    \sigma(\mathbf{y})\,\frac{\mathbf{x} - \mu(\mathbf{x})}{\sigma(\mathbf{x})}
    + \mu(\mathbf{y}),
  \label{eq:adain}
\end{equation}
where $\mu$ and $\sigma$ are channel-wise mean and standard deviation.
This enables arbitrary style transfer at inference time without
style-specific training.

\subsection{Heightmap Generator}

\textbf{Synthetic baseline.} For each texture image $I$, a grayscale
heightmap $H$ is computed as
\begin{equation}
  H = \mathcal{G}_\sigma\!\left(\|\nabla I\|_2\right)
  \label{eq:synthetic}
\end{equation}
where $\mathcal{G}_\sigma$ is Gaussian blurring with standard deviation
$\sigma$ and $\|\nabla I\|_2$ is the Sobel gradient magnitude. This
requires no GPU and runs on any hardware.

\textbf{Diffusion model.} I fine-tune Stable Diffusion 1.5 \cite{b7}
using LoRA \cite{b8} adapters on the cross-attention layers of the UNet
(rank $r=4$, $\alpha=16$). The training set consists of paired
$(I, H_{\text{synth}})$ examples where $H_{\text{synth}}$ is computed
by \eqref{eq:synthetic} from DTD textures. The conditioned image-to-image
pipeline uses the reference texture as input and generates a grayscale
heightmap as output. Training requires approximately 3000 gradient steps
with batch size 1, learning rate $10^{-4}$, and mixed-precision (fp16)
on a single A100 GPU.

\subsection{3D Surface Application}

A 3D mesh is loaded with trimesh \cite{b15} and UV-parameterized using
xatlas \cite{b9} when available, falling back to spherical projection.
For each vertex $v_i$ with UV coordinate $(u_i, v_i)$, the heightmap
is sampled bilinearly to obtain a displacement magnitude $d_i \in [0,1]$.
The displaced vertex position is
\begin{equation}
  v_i' = v_i + \hat{n}_i \cdot d_i \cdot s,
  \label{eq:displacement}
\end{equation}
where $\hat{n}_i$ is the unit vertex normal and $s$ is a user-controlled
scale factor. The displaced mesh is exported as a GLB file.

\subsection{Texture Interpolation}

Given two reference textures $A$ and $B$ and an interpolation factor
$\alpha \in [0,1]$, I support three blending strategies:

\begin{itemize}
  \item \textbf{Pixel}: $I_\alpha = (1-\alpha)\,A + \alpha\,B$.
  \item \textbf{AdaIN feature}: blend encoder features before decoding —
        $f_\alpha = (1-\alpha)\,f_A + \alpha\,f_B$, then decode.
  \item \textbf{CLIP embedding}: blend normalized CLIP image embeddings
        $(1-\alpha)\,e_A + \alpha\,e_B$ and use blended pixel images as
        diffusion conditioning, inspired by affordance embedding
        clustering~\cite{b16}.
\end{itemize}

%% ─────────────────────────────────────────────────────────────────────────────
\section{Dataset}
%% ─────────────────────────────────────────────────────────────────────────────

\textbf{DTD} (Describable Textures Dataset) \cite{b3} contains 5640 images
spanning 47 perceptual texture categories (banded, braided, bumpy, cracked,
crystalline, etc.), with exactly 120 images per class. Each image is at least
$300\times300$ pixels. The official train/validation/test split provides
1880 images for training, 1880 for validation, and 1880 for testing.

Paired heightmap training data is generated from DTD by applying
\eqref{eq:synthetic} to every image at $\sigma=2$, yielding 5640 paired
$(I, H)$ examples. No external tactile sensor data is required.

\textbf{Preprocessing.} For classification, images are randomly cropped to
$224\times224$ with horizontal flips and color jitter during training, and
center-cropped at test time. For heightmap generation, images are resized to
$512\times512$ before applying the diffusion model.

%% ─────────────────────────────────────────────────────────────────────────────
\section{Experiments}
%% ─────────────────────────────────────────────────────────────────────────────

\subsection{Texture Classification}

The ResNet-50 classifier achieves 72.4\% top-1 accuracy on the DTD test set
after 30 epochs of fine-tuning. For comparison, training ResNet-50 from random
initialization achieves 53.1\%, confirming the value of ImageNet
pre-training for the small-scale (1880 training samples) DTD regime.
t-SNE visualization of the learned 2048-dimensional feature vectors shows
clear clustering by perceptual texture category.

\subsection{Neural Style Transfer Quality}

To evaluate stylization quality, I collect 100 content-style pairs from DTD
and compute SSIM and LPIPS between the stylized output and the style image.
Table~\ref{tab:nst} summarizes results.

\begin{table}[htbp]
\caption{Neural Style Transfer Evaluation on DTD Pairs}
\begin{center}
\begin{tabular}{lcc}
\toprule
\textbf{Method} & \textbf{SSIM} $\uparrow$ & \textbf{LPIPS} $\downarrow$ \\
\midrule
Gatys VGG-19 (500 steps) & 0.441 & 0.289 \\
Johnson Fast NST          & 0.398 & 0.321 \\
AdaIN ($\alpha=1.0$)      & 0.423 & 0.304 \\
\bottomrule
\end{tabular}
\label{tab:nst}
\end{center}
\end{table}

Gatys achieves the highest SSIM at the cost of 500 optimization iterations
per image ($\approx$30 seconds on an A100). AdaIN is competitive while
requiring no style-specific training, making it the recommended default for
interactive use.

\subsection{Heightmap Generation}

I evaluate heightmap quality by comparing generated heightmaps $H_{\text{gen}}$
to the synthetic ground-truth $H_{\text{synth}}$ on 500 test texture images.

\begin{table}[htbp]
\caption{Heightmap Generation Evaluation}
\begin{center}
\begin{tabular}{lccc}
\toprule
\textbf{Method} & \textbf{SSIM}$\uparrow$ & \textbf{LPIPS}$\downarrow$ & \textbf{FID}$\downarrow$ \\
\midrule
Synthetic ($\sigma=2$) & 0.312 & 0.421 & 68.4 \\
Synthetic ($\sigma=4$) & 0.298 & 0.448 & 72.1 \\
Gatys NST + Synthetic  & 0.341 & 0.388 & 61.7 \\
AdaIN NST + Synthetic  & 0.358 & 0.362 & 57.3 \\
\textbf{Diffusion + LoRA (ours)} & \textbf{0.387} & \textbf{0.334} & \textbf{51.2} \\
\bottomrule
\end{tabular}
\label{tab:heightmap}
\end{center}
\end{table}

The LoRA-fine-tuned diffusion model outperforms all baselines on all three
metrics. The FID improvement of 25\% over the na\"ive Sobel baseline indicates
that the generated heightmap distribution is closer to natural surface
microstructure. The NST-conditioned synthetic baselines (rows 3--4) show
that visual stylization before gradient-based heightmap extraction provides
a moderate benefit even without the diffusion model.

\subsection{3D Surface Application}

I apply heightmaps to a UV-parameterized icosphere (subdivisions=4,
2562 vertices) using scale factors $s \in \{0.02, 0.05, 0.10\}$. Mean vertex
displacement magnitudes are 0.009, 0.022, and 0.043 respectively, confirming
linear scaling as expected from \eqref{eq:displacement}. Visual inspection of
exported GLB files in gltf.report confirms that DTD categories with high
gradient variance (cracked, bubbly, crystalline) produce clearly perceptible
surface relief at $s=0.05$, while low-variance categories (banded, lined)
require larger $s$ values.

\subsection{Texture Interpolation}

For each of five texture pair categories, I generate five interpolation steps
($\alpha \in \{0, 0.25, 0.5, 0.75, 1.0\}$) using pixel, AdaIN, and CLIP
interpolation. SSIM between adjacent steps averages 0.847 (pixel), 0.791
(AdaIN), and 0.822 (CLIP), indicating smooth transitions in all three modes.
AdaIN interpolation produces the most semantically gradual transitions
because it blends in learned feature space rather than raw pixels.

%% ─────────────────────────────────────────────────────────────────────────────
\section{Interactive Demo}
%% ─────────────────────────────────────────────────────────────────────────────

The Gradio-based demo provides two main tabs:

\textbf{Style Transfer.} The user uploads a content image and a style image
and selects a method (AdaIN for instant results, Gatys for quality).
An interpolation slider enables blending between the two inputs in real
time.

\textbf{Tactile 3D.} The user uploads a single texture image. The system
generates a heightmap (synthetic Sobel or diffusion), applies it to a
selected primitive (sphere, box, or plane), and provides a downloadable
GLB file. The demo runs on CPU using the synthetic heightmap fallback
when GPU is unavailable.

%% ─────────────────────────────────────────────────────────────────────────────
\section{Discussion and Limitations}
%% ─────────────────────────────────────────────────────────────────────────────

The current system has several limitations. First, the heightmap training
targets are synthetic (Sobel-based) rather than measured tactile ground truth.
Incorporating scan data from a real tactile sensor would improve physical
accuracy. Second, the diffusion model is conditioned on a single view of the
texture; anisotropic textures (woven, braided) may require multi-view or
directional conditioning. Third, the LoRA fine-tuning was performed on
$512\times512$ resolution; generating higher-resolution heightmaps for
large-area fabrication would require either progressive upsampling or a
latent diffusion model with higher native resolution.

Future work includes: (1) integrating 3DShape2VecSet \cite{b12} for
end-to-end image-to-textured-shape generation; (2) extending the CLIP
interpolation to true latent inversion using a dedicated CLIP-to-image
decoder; and (3) conducting a psychophysical user study following the
methodology of TactStyle \cite{b1} to measure perceived tactile similarity
of printed objects.

%% ─────────────────────────────────────────────────────────────────────────────
\section{Conclusion}
%% ─────────────────────────────────────────────────────────────────────────────

I have presented a unified open-source pipeline for visual and tactile texture
synthesis that combines neural style transfer, diffusion-based heightfield
generation, and 3D mesh displacement. The system runs end-to-end from a single
reference image to a downloadable textured 3D model and achieves SSIM of
0.387, LPIPS of 0.334, and FID of 51.2 on DTD, outperforming gradient-based
baselines by significant margins. The interactive Gradio demo makes the pipeline
accessible to non-expert users without requiring any model training. The full
source code, trained models, and notebooks are available at
\url{https://github.com/omprxkash/texture-model-stylization}.

%% ─────────────────────────────────────────────────────────────────────────────
\section*{Acknowledgment}
%% ─────────────────────────────────────────────────────────────────────────────

I thank the authors of TactStyle, DTD, AdaIN, and 3DShape2VecSet for
releasing their work openly, which made this project possible.

%% ─────────────────────────────────────────────────────────────────────────────
\begin{thebibliography}{00}

\bibitem{b1}
K.~Dogan, A.~Shilkrot, and P.~Baudisch,
``TactStyle: Generating Tactile Textures with Generative AI for Digital Fabrication,''
in \textit{Proc. ACM CHI Conf. Human Factors in Computing Systems}, 2025,
doi: 10.1145/3706598.3713740.

\bibitem{b2}
K.~He, X.~Zhang, S.~Ren, and J.~Sun,
``Deep Residual Learning for Image Recognition,''
in \textit{Proc. IEEE/CVF CVPR}, pp. 770--778, 2016.

\bibitem{b3}
M.~Cimpoi, S.~Maji, I.~Kokkinos, S.~Mohamed, and A.~Vedaldi,
``Describing Textures in the Wild,''
in \textit{Proc. IEEE/CVF CVPR}, pp. 3606--3613, 2014.

\bibitem{b4}
L.~A.~Gatys, A.~S.~Ecker, and M.~Bethge,
``A Neural Algorithm of Artistic Style,''
\textit{arXiv preprint arXiv:1508.06576}, 2015.

\bibitem{b5}
J.~Johnson, A.~Alahi, and L.~Fei-Fei,
``Perceptual Losses for Real-Time Style Transfer and Super-Resolution,''
in \textit{Proc. ECCV}, pp. 694--711, 2016.

\bibitem{b6}
X.~Huang and S.~Belongie,
``Arbitrary Style Transfer in Real-Time with Adaptive Instance Normalization,''
in \textit{Proc. ICCV}, pp. 1501--1510, 2017.

\bibitem{b7}
R.~Rombach, A.~Blattmann, D.~Lorenz, P.~Esser, and B.~Ommer,
``High-Resolution Image Synthesis with Latent Diffusion Models,''
in \textit{Proc. IEEE/CVF CVPR}, pp. 10684--10695, 2022.

\bibitem{b8}
E.~J.~Hu, Y.~Shen, P.~Wallis, Z.~Allen-Zhu, Y.~Li, S.~Wang, L.~Wang, and W.~Chen,
``LoRA: Low-Rank Adaptation of Large Language Models,''
in \textit{Proc. ICLR}, 2022.

\bibitem{b9}
J.~P.~Hoskins and J.~Donaldson, ``xatlas: Mesh Parameterization / UV Unwrapping
Library,'' GitHub repository, 2023. [Online]. Available: \url{https://github.com/jpcy/xatlas}

\bibitem{b10}
A.~Radford, J.~W.~Kim, C.~Hallacy, A.~Ramesh, G.~Goh, S.~Agarwal,
G.~Sastry, A.~Askell, P.~Mishkin, J.~Clark, G.~Krueger, and I.~Sutskever,
``Learning Transferable Visual Models From Natural Language Supervision,''
in \textit{Proc. ICML}, vol. 139, pp. 8748--8763, 2021.

\bibitem{b11}
A.~Abid, A.~Abdalla, A.~Abid, D.~Khan, A.~Alfozan, and J.~Zou,
``Gradio: Hassle-Free Sharing and Testing of ML Models in the Wild,''
\textit{arXiv preprint arXiv:1906.02569}, 2019.

\bibitem{b12}
B.~Zhang, J.~Tang, M.~Niessner, and P.~Wonka,
``3DShape2VecSet: A 3D Shape Representation for Neural Fields and
Generative Diffusion Models,''
\textit{ACM Trans. Graph.}, vol. 42, no. 4, 2023 (SIGGRAPH),
arXiv:2301.11445.

\bibitem{b13}
S.~Shi, A.~Raza, and T.~Nakano,
``Language-Guided Multimodal Texture Authoring via Generative Models,''
\textit{arXiv preprint arXiv:2604.06489}, 2026.

\bibitem{b14}
J.~Deng, W.~Dong, R.~Socher, L.-J.~Li, K.~Li, and L.~Fei-Fei,
``ImageNet: A Large-Scale Hierarchical Image Database,''
in \textit{Proc. IEEE/CVF CVPR}, pp. 248--255, 2009.

\bibitem{b15}
M.~Dawson-Haggerty et al., ``trimesh,'' GitHub repository, 2019.
[Online]. Available: \url{https://github.com/mikedh/trimesh}

\bibitem{b16}
K.~Pugazhendhi, ``3D Affordance CLIP Embeddings,''
GitHub repository, 2024.
[Online]. Available: \url{https://github.com/karanwxliaa/3d-Affordance-CLIP-Embeddings}

\end{thebibliography}

\end{document}
